\section{Notation and system model}

Let $G=(V,E)$ be a finite undirected graph, with $n=|V|$ vertices and
$m=|E|$ edges. Let $C=\{1,2,\ldots,k\}$ be a finite color set. A coloring, or
configuration, is denoted by
\[
    c=(c_i)_{i\in V}\in C^V,
\]
where $c_i$ is the color assigned to vertex $i$. For a vertex $i\in V$, let
$N(i)=\{j\in V:\{i,j\}\in E\}$ be its neighborhood. The local conflict cost of
vertex $i$ under configuration $c$ is
\[
    d_i(c_i,c_{-i})
    =
    \sum_{j\in N(i)} \mathbf{1}\{c_i=c_j\},
\]
where $c_{-i}$ denotes the colors of all vertices except $i$. When the colors
of the other vertices are fixed, we also write
\[
    d_i(\gamma,c_{-i})
    =
    \sum_{j\in N(i)} \mathbf{1}\{\gamma=c_j\},
    \qquad \gamma\in C .
\]

The competitive recoloring dynamics are asynchronous. At each activation, a
single vertex $i$ is selected. If there exists a color that strictly reduces
its local conflict cost, vertex $i$ changes to a best-response color:
\[
    c_i^+
    \in
    \argmin_{\gamma\in C} d_i(\gamma,c_{-i}).
\]
The update is accepted only when
\[
    d_i(c_i^+,c_{-i}) < d_i(c_i,c_{-i});
\]
otherwise the configuration is left unchanged. A \emph{step} below means one
accepted recoloring move.

Define the global conflict potential by
\[
    \Phi(c)
    =
    \sum_{\{u,v\}\in E} \mathbf{1}\{c_u=c_v\}
    =
    \frac12 \sum_{i\in V} d_i(c_i,c_{-i}).
\]
Thus $\Phi(c)$ is exactly the number of monochromatic edges in $G$.

The global minimization of $\Phi$ is computationally nontrivial: for
$k\ge 2$, minimizing monochromatic edges is equivalent to maximizing the
number of bichromatic edges, i.e., to the MAX-$k$-CUT objective. Since this
problem is NP-hard in general, a distributed best-response rule that is
guaranteed to reach a locally stable configuration in finite time is valuable
even though it does not claim global optimality.

\begin{lemma}[Exact potential drop]
Let $c$ and $c^+$ be two configurations that differ only at one vertex $i$.
Suppose vertex $i$ changes from color $\alpha=c_i$ to color $\beta=c_i^+$ and
strictly reduces its local conflict cost, while $c_{-i}^+=c_{-i}$. Then
\[
    \Phi(c^+)-\Phi(c)
    =
    d_i(\beta,c_{-i})-d_i(\alpha,c_{-i})
    \le -1 .
\]
\end{lemma}

\begin{proof}
Only edges incident to $i$ can change their conflict status, since all other
vertices retain their colors. Hence every edge not in
$\{\{i,j\}:j\in N(i)\}$ contributes the same amount to $\Phi(c^+)$ and
$\Phi(c)$.

Before the update, the number of conflicting edges incident to $i$ is
\[
    d_i(\alpha,c_{-i})
    =
    \sum_{j\in N(i)} \mathbf{1}\{\alpha=c_j\}.
\]
After the update, the number of conflicting edges incident to $i$ is
\[
    d_i(\beta,c_{-i})
    =
    \sum_{j\in N(i)} \mathbf{1}\{\beta=c_j\}.
\]
Therefore the total change in the number of monochromatic edges is exactly
\[
    \Phi(c^+)-\Phi(c)
    =
    d_i(\beta,c_{-i})-d_i(\alpha,c_{-i}).
\]
By the strict-improvement rule,
$d_i(\beta,c_{-i})<d_i(\alpha,c_{-i})$. Since both quantities are integers,
their difference is at most $-1$. This proves the claim.
\end{proof}

\begin{theorem}[Finite convergence]
For any finite graph $G=(V,E)$, any finite color set $C$, and any initial
configuration $c(0)\in C^V$, the competitive recoloring dynamics terminate
after at most $|E|$ accepted recoloring steps. More precisely, if $T$ denotes
the number of accepted recolorings before termination, then
\[
    T \le \Phi(c(0)) \le |E|.
\]
The terminal configuration is a pure Nash equilibrium, equivalently a fixed
point of the strict best-response dynamics.
\end{theorem}

\begin{proof}
By Lemma~1, every accepted recoloring decreases $\Phi$ by at least one. The
potential is integer-valued and nonnegative because it counts monochromatic
edges:
\[
    0 \le \Phi(c) \le |E|
    \qquad \text{for every } c\in C^V .
\]
Thus a sequence of accepted recolorings can have length no larger than the
initial potential value. Consequently,
\[
    T \le \Phi(c(0)) \le |E|.
\]
In particular, no infinite sequence of accepted recolorings is possible, so
the dynamics terminate in finite time.

At termination, no vertex has a color in $C$ that strictly lowers its local
conflict cost. Hence no individual vertex can improve by a unilateral
recoloring while the other vertices remain fixed. This is precisely a pure
Nash equilibrium for the local conflict-cost game, and equivalently a fixed
point of the stated strict best-response dynamics.
\end{proof}

\begin{remark}[Average-case scaling and empirical validation]
Theorem~1 gives a deterministic worst-case upper bound of $|E|$ accepted
recoloring steps. The accompanying simulations examine randomized asynchronous
updates on sparse Erd\H{o}s--Renyi and Barabasi--Albert graphs with constant
average degree, for which $|E|=\mathcal{O}(|V|)$. Across the reported trials,
the observed convergence time is substantially below the worst-case edge
bound, typically about $0.12|E|$ accepted recolorings. These data support the
interpretation that, under randomized asynchronous updates, the adversarial
activation orders needed to approach the worst-case bound are rarely sampled in
standard sparse graph ensembles. Consequently, the empirical behavior is closer
to $\mathcal{O}(|V|\log |V|)$, and in these experiments nearly linear in the
network size, although the formal guarantee used here remains the deterministic
edge bound above.
\end{remark}
